% ver3 (2026-11-17 night): English version of ja_v3/09_結論.md (more specific counts than the abstract, after reviews/v2_ja/prose_body_okada_ishii D; ends with the proposal's result).
\section{Conclusion}\label{sec:conclusion}

We treated multimodal emotion recognition on a new corpus with a few labels per emotion as semi-supervised domain adaptation and answered three questions with linear classifiers on frozen pretrained encoders. With one source corpus, audio and visual answer bias was almost removed by subtracting feature means without labels; for audio at K = 5, mean subtraction gained 4.69\pt, close to the 5.26\pt\ of the merged classifier. With all other groups as sources, conditions in which the remaining error was reduced numbered 7 of 210 for eight methods using same-modality unlabelled utterances and 38 of 81 for pseudo-labels from classifiers including other modalities. With AFEW, IEMOCAP (improvised), MOSEI and MELD each as the new corpus, distillation on these pseudo-labels gave the audio classifier 1.4--4.1\pt\ higher UAR than before retraining in all 12 conditions with 5, 10 and 25 labels per emotion (point estimates), with intervals above zero in the nine conditions other than IEMOCAP. Against 11 methods that do not use pseudo-labels including other modalities, distillation was higher in all 132 comparisons in point estimates and in 122 with intervals above zero. Across 360 post hoc pseudo-label variants per K (K = 5, 25), the rank correlation between the advantage of the pseudo-labels and the UAR difference was 0.80 and 0.83, and no variant whose pseudo-labels had lower UAR than the classifier before retraining (173 and 184) rose. Finally, because the strongest modality differed by corpus and emotion, we proposed late fusion that measures strengths on the K target labels and per-emotion profiles on the sources, with hyperparameters selected on source-corpus episodes. The AVT retrained on its pseudo-labels had higher UAR than the AVT retrained on the product in 11 of 12 conditions in point estimates. The estimation-time UAR of the proposal was higher than that of the product of the three classifiers' probabilities in 11 of 12 conditions in point estimates, with intervals excluding zero in seven, and by 2.1--3.5\pt\ on AFEW.
